\documentclass[a4paper,12pt]{article}
\usepackage[utf8]{inputenc}
\usepackage[T1]{fontenc}
\usepackage[ngerman]{babel}
\usepackage{geometry}
\usepackage{amsmath}
\usepackage{amssymb}
\usepackage{booktabs}
\geometry{margin=2.5cm}

\begin{document}

% ================================================================
\section*{GNOT: A General Neural Operator Transformer for Operator Learning}

\textbf{Autoren:} Zhongkai Hao, Zhengyi Wang, Hang Su, Chengyang Ying, Yinpeng Dong, \\
\hspace*{2.3cm} Songming Liu, Ze Cheng, Jian Song, Jun Zhu \\
\textbf{Institution:} Tsinghua-Bosch Joint ML Center, Tsinghua University; RealAI; Bosch China Investment Ltd. \\
\textbf{Konferenz:} ICML 2023 \\
\textbf{arXiv:} 2302.14376

% ================================================================
\subsection*{1\quad Problemstellung und Motivation}

Bestehende Operatorlerner für PDEs weisen drei praktische Einschränkungen auf:
\begin{enumerate}
    \item \textbf{Irreguläre Gitter:} Methoden wie FNO und U-Net sind auf reguläre
    rechteckige Gitter beschränkt. Reale Geometrien (z.B.\ Tragflügelprofile, Wärmesenken)
    erfordern unstrukturierte Gitter, auf denen diese Ansätze versagen.

    \item \textbf{Multiple Eingabefunktionen:} Physikalische Systeme werden oft durch
    mehrere Eingabegrößen beschrieben (Randform, verteilte Quellterme, Materialparameter).
    Bestehende Methoden wie FNO und MIONet sind nicht direkt auf diese allgemeine
    Situation ausgelegt.

    \item \textbf{Mehrskaligkeitsprobleme:} In vielen Systemen sind unterschiedliche
    physikalische Phänomene auf verschiedenen Skalen aktiv (z.B.\ Nahfeld vs.\ Fernfeld
    bei Strömungen). Ein einzelnes globales Modell kann solche Strukturen nur schwer
    gleichzeitig repräsentieren.
\end{enumerate}

GNOT (\emph{General Neural Operator Transformer}) adressiert alle drei Herausforderungen
durch zwei neuartige Komponenten: (1)~einen \emph{Heterogeneous Normalized Attention}-Block
(HNA) für allgemeines Operatorlernen mit variabler Anzahl von Eingabefunktionen und
(2)~einen \emph{Geometric Gating}-Mechanismus (GGM) für mehrskalige Probleme.

% ================================================================
\subsection*{2\quad Methodik: HNA-Block und geometrisches Gating}

\paragraph{Eingabekodierung.}
Das Modell empfängt Abfragepunkte $\{x_i^q\}_{i=1}^{N_q}$ und bis zu $L$ verschiedene
Eingabefunktionen $\{a^1, \ldots, a^L\}$, die auf potenziell unterschiedlichen Gittern
$\{x^{j}\}$ diskretisiert sind. Jede Eingabefunktion wird durch ein individuelles
MLP $f_w(x_i, a_i)$ in ein Embedding $Y^l \in \mathbb{R}^{N_l \times n_e}$ kodiert.

\paragraph{Heterogeneous Normalized Cross-Attention.}
Der HNA-Block berechnet zunächst eine normalisierte lineare Cross-Attention zwischen
den Abfragepunkten $X \in \mathbb{R}^{N_q \times n_e}$ und den $L$ Konditionierungs-Embeddings
$\{Y^l\}$:
\begin{align}
    \bar{q}_i &= \mathrm{Softmax}(q_i) = \left(\frac{e^{q_{ij}}}{\sum_j e^{q_{ij}}}\right),\quad
    \bar{k}_i = \mathrm{Softmax}(k_i), \\
    z_t &= \bar{q}_t + \frac{1}{L} \sum_{l=1}^{L} \alpha_t^l\,\bar{q}_t \cdot
           \left(\sum_{i_l=1}^{N_l} \bar{k}_{i_l}^l \otimes v_{i_l}^l\right),
\end{align}
wobei $\alpha_t^l = \left(\sum_j \bar{q}_t \cdot \bar{k}_j^l\right)^{-1}$ der
Normierungskoeffizient ist. Der Aufwand ist $\mathcal{O}((N_q + \sum_l N_l) n_e^2)$,
also linear in der Sequenzlänge. Anschließend wird eine normalisierte lineare
Selbst-Attention auf $X$ angewendet. Dieser \emph{Cross+Self}-Block ist der
leistungsstärkste Aufmerksamkeitstyp im Ablationsvergleich.

\paragraph{Geometrisches Gating.}
Um mehrskalige Probleme zu behandeln, wird ein \emph{Mixture-of-Experts}-Mechanismus
mit geometrischer Konditionierung eingesetzt. In jedem Block gibt es $K$ spezialisierte
FFN-Expertennetze $\{E_i(\cdot)\}$; ein Gating-Netzwerk $G(\cdot) : \mathbb{R}^d \to \mathbb{R}^K$
berechnet die räumlich variablen Gewichte:
\begin{equation}
    z_t \leftarrow z_t + \sum_{i=1}^{K} p_i(x_t) \cdot E_i(z_t), \quad
    p_i(x_t) = \frac{\exp(G_i(x_t))}{\sum_{j=1}^K \exp(G_j(x_t))}.
\end{equation}
Das Gating-Netzwerk erhält die geometrischen Koordinaten der Abfragepunkte als Eingabe
und lernt so eine weiche Gebietsdekomposition, bei der unterschiedliche Skalen
durch spezialisierte Experten abgedeckt werden.

% ================================================================
\subsection*{3\quad Experimente und Ergebnisse}

GNOT wird auf neun Datensätzen aus vier Bereichen (Strömungsmechanik, Elastizität,
Elektromagnetismus, Wärmeleitung) getestet:

\begin{center}
\begin{tabular}{lrrrr}
\toprule
\textbf{Datensatz} & \textbf{MIONet} & \textbf{FNO} & \textbf{GK-Transf.} & \textbf{GNOT} \\
\midrule
Darcy2d & 5.45e-2 & 1.09e-2 & \textbf{8.40e-3} & 1.05e-2 \\
NS2d (voll) & -- & 8.20e-2 & 7.92e-2 & \textbf{4.43e-2} \\
Elasticity & 9.65e-2 & 5.08e-2 & 2.01e-2 & \textbf{8.65e-3} \\
NACA & 1.32e-1 & 4.21e-2 & 1.61e-2 & \textbf{7.57e-3} \\
Inductor2d ($A_z$) & 3.10e-2 & -- & 2.56e-1 & \textbf{1.21e-2} \\
Heat (voll) & 1.45e-1 & -- & -- & \textbf{2.56e-2} \\
\bottomrule
\end{tabular}
\end{center}

Für Datensätze mit irregulärem Gitter und multiplen Eingaben (Typen A und/oder B)
reduziert GNOT den Vorhersagefehler gegenüber den Basislinien um typischerweise
40--50\,\%. Auf dem einfachen regulären Darcy2d-Gitter ist der Galerkin-Transformer
minimal besser. Bei mehrskaligen Problemen (Heatsink, Heat) kann GNOT als einzige
Methode überhaupt eine sinnvolle Lösung liefern.

Skalierungsexperimente zeigen, dass GNOT data-effizienter als MIONet ist:
Mit wachsender Datenmenge nähert sich der Fehler schneller dem Optimum, und
größere Modellkapazität (GNOT-large) verbessert die Leistung konsistent gemäß
einer annähernd logarithmischen Skalierungsregel.

% ================================================================
\subsection*{4\quad Bezug zur Masterarbeit}

\paragraph{4.1\quad Multiple Eingabekanäle als konzeptuelle Parallele.}
In der Masterarbeit wird das Strömungsfeld aus fünf Eingabekanälen (Kristallmaske,
Signed-Distance-Feld, Abstandsfeld zum nächsten Kristall, normalisierte $x$- und $z$-Koordinaten)
vorhergesagt. Dies entspricht einer mehrkanaligen Eingabestruktur, die konzeptuell
mit GNOT's Fähigkeit zur Verarbeitung multipler Eingabefunktionen verwandt ist.
Der HNA-Block bietet dabei den Vorteil, dass jeder Eingabekanal durch ein separates
Encoder-MLP verarbeitet wird und die Kreuzaufmerksamkeit die Beiträge flexibel
gewichtet -- eine Eigenschaft, die für die Modellierung komplexer Wechselwirkungen
zwischen Kristallgeometrie und Strömungsfeld interessant sein könnte.

\paragraph{4.2\quad Irreguläre Gitter und variierende Topologie.}
Ein zentrales Ziel der Masterarbeit ist die Generalisierung auf variierende Kristallzahlen,
also auf Konfigurationen mit unterschiedlicher geometrischer Topologie. GNOT ist so
konzipiert, dass es auf irregulären Gittern mit beliebigen Abfragepunkten operiert.
Diese Eigenschaft würde es prinzipiell erlauben, Kristallkonfigurationen direkt durch
ihre geometrischen Gitterpunkte (z.B.\ Knoten eines unstrukturierten FEM-Netzes) zu
repräsentieren -- eine Alternative zur regulären Gitterrepräsentation der Masterarbeit.
Allerdings setzt dies eine deutlich andere Datenpipeline voraus und wurde in der
Masterarbeit bewusst zugunsten der regulären $256 \times 256$-Repräsentation verzichtet.

\paragraph{4.3\quad Geometrisches Gating als Mehrskaligkeitsstrategie.}
In der Multi-Kristall-Stokes-Strömung koexistieren zwei charakteristische Längenskalen:
die lokalen Grenzschichten unmittelbar an den Kristalloberflächen und die
langreichweitigen Wirbelstrukturen im Fernfeld. Der geometrische Gating-Mechanismus
in GNOT -- der durch räumliche Koordinaten eine weiche Gebietsdekomposition lernt --
ist strukturell geeignet, diese Mehrskaligkeit zu adressieren: Verschiedene Experten
könnten sich auf Nah- und Fernfeldregionen spezialisieren. Dies ist eine konzeptuell
interessante Alternative zur Mehrskalenverarbeitung durch Skip-Connections im U-Net.

\paragraph{4.4\quad Skalierungsverhalten und Dateneffizienz.}
Die in der Masterarbeit beobachteten Unterschiede zwischen U-Net und FNO in Bezug auf
Generalisierungsverhalten könnten durch einen Architekturvergleich mit GNOT ergänzt
werden. Das Skalierungsverhalten von GNOT -- bessere Nutzung größerer Datensätze bei
höherer Modellkapazität -- könnte relevant sein, wenn in zukünftigen Arbeiten die
Trainingsdatengröße erhöht wird.

\paragraph{4.5\quad Abgrenzung.}
GNOT wurde für Probleme mit irregulären Gittern und multiplen Eingabefunktionen auf
\emph{fixer} Geometrie entwickelt. Die variierende Topologie (unterschiedliche Anzahl
eingebetteter Festkörper) ist ein anderes Generalisierungsproblem als die in GNOT
adressierte Skalierungsheterogenität. Insbesondere enthält GNOT keine Mechanismen
zur expliziten Codierung der Anzahl oder Positionen von eingebetteten Festkörpern --
ein Problem, das in der Masterarbeit durch Kristallmasken und Abstandsfelder gelöst wird.
Die Frage, ob GNOTs allgemeines Encoder-Framework auf variierende Kristallzahlen
generalisiert, wäre ein eigenständiges Forschungsprojekt.

\end{document}
